\documentclass[11.5pt]{article}

\usepackage{amsmath, amsthm, amssymb, amsfonts}
\usepackage{graphicx}
\usepackage{setspace}
\usepackage{geometry}
\usepackage{float}
\usepackage{hyperref}
\usepackage[utf8]{inputenc}
\usepackage[english]{babel}
\usepackage[dvipsnames]{xcolor}
\usepackage{xcolor}

\DeclareFixedFont{\ttb}{T1}{txtt}{bx}{n}{12}
\DeclareFixedFont{\ttm}{T1}{txtt}{m}{n}{12}

\hypersetup { colorlinks, linkcolor=Violet, urlcolor=Cyan}

\usepackage{color}
\definecolor{deepblue}{rgb}{0,0,0.5}
\definecolor{deepred}{rgb}{0.6,0,0}
\definecolor{deepgreen}{rgb}{0,0.5,0}

\usepackage{listings}
\usepackage{color}
\definecolor{dkgreen}{rgb}{0,0.6,0}
\definecolor{gray}{rgb}{0.5,0.5,0.5}
\definecolor{mauve}{rgb}{0.58,0,0.82}

\colorlet{LightGreen}{Green!15}

\newcommand{\HRule}[1]{\rule{\linewidth}{#1}}

\setstretch{1.2}
\geometry{
	textheight=9in,
	textwidth=5.5in,
	top=1in,
	headheight=12pt,
	headsep=25pt,
	footskip=30pt
}

% ------------------------------------------------------------------------------

\begin{document}
	
	\pagenumbering{gobble}
	
	
	\title{\includegraphics[height=4.5cm]{IUT_logo}
		\normalsize \textsc{}
		\\ [1.5cm]
		\HRule{1.5pt} \\
		
		
		\LARGE \textbf{\uppercase{IUT-Medical Management System} MediEase
			\HRule{1.5pt} \\ [1.3cm] \LARGE{SWE 4404 \\ Software Project Lab II} \vspace*{1\baselineskip}}
	}
	
	\date{\LARGE\today}
	
	\author{\HRule{1.5pt} \\
		\\
		\\\Large\textbf{{Ahmed Alfey Sani} [210042129]}\\
		\Large\textbf{{Shefayat E - Shams Adib} [210042141]}\\
		\Large\textbf{{Hasin Mahtab Alvee} [210042174]}\\
		\\
		\Large\textbf{Supervisor: Imtiaj Ahmed Chowdhury}
		\\
		\\
		\LARGE Department of CSE \\
		\LARGE B.Sc in Software Engineering
	}
	
	
	\maketitle
	\newpage
	\tableofcontents
	\newpage
	\pagenumbering{arabic}
	% ------------------------------------------------------------------------------
	
	\section{Project Overview}
	
	IUT-Medical Management System is a comprehensive solution designed to streamline operations and enhance patient care within the Islamic University of Technology's medical center. It offers features such as appointment scheduling, medical expense management, room reservations, prescription generation for doctors, and availability tracking for doctors.
	
	
	
	\section{Motivation behind the Project}
	
	The motivation behind IUT Medical Center Digital Platform project stems from a commitment to modernize and optimize healthcare delivery at the Islamic University of Technology's medical center. Recognizing the evolving landscape of healthcare technology and the increasing demand for streamlined patient services, this project aims to address key challenges faced by both patients and healthcare providers. By implementing a comprehensive digital solution, we seek to enhance accessibility, convenience, and efficiency in accessing medical care while also improving operational workflows and resource management within the medical center. Through this initiative, we aspire to elevate the overall patient experience, foster stronger doctor-patient relationships, and ultimately contribute to the delivery of high-quality healthcare services that meet the evolving needs of our community.

	
	% ------------------------------------------------------------------------------
	
	% ------------------------------------------------------------------------------
	
	
	\section{Key Features}
	
		
	\subsection{Generate Prescription}
	
	Doctors can prescribe the necessary medicines for patients and generate digital prescriptions directly within the platform from the medicine database which will include a variety of medicines based on different types and needs.
	
	\subsection{Manage Medical Records}
	
	This feature empowers doctors and nurses to efficiently manage and update medical records for patients under their care within the IUT Medical Center Digital Platform. It provides nurses with secure access to patient information, streamlining documentation processes and ensuring accurate and up-to-date medical records.
	
	\subsection{Due Medical Expense Management}

	The Due Medical Expense Management feature enables students with referrals to submit medical bills incurred outside IUT for reimbursement. Students upload referral documents, prescriptions, and bills for administrative verification. Once verified, administrators clear dues. This feature serves students and admins, updating due trackers in student profiles for transparency.
	
	
	\subsection{Appointment Scheduler}
	
	Patients can conveniently book appointments with healthcare providers through an intuitive online calender. The system provides real-time availability of doctors, allowing patients to select preferred time slots based on their convenience.
	
	
	\subsection{Service Reservations}
	
	Patients can reserve rooms or bed or other available services through the platform. Hospital staff can also efficiently manage reservations and services for patients.
	
	
	\subsection{Staff Scheduler}
	This feature enables efficient scheduling and management of healthcare staff within the IUT Medical Center. It provides administrators with tools to create, view, and adjust staff schedules, ensuring adequate coverage and optimal resource allocation across various departments and shifts.
	
	
	% ------------------------------------------------------------------------------
	
	% ------------------------------------------------------------------------------
	
	\section{Tools and Technology}
	
	
	\subsection{Frontend}
	We will be using \textbf{ReactJS} as the primary Javascript framework for our frontend technology along with TailwindCSS for styling the client side.
	
	\subsection{Backend}
	We will be using \textbf{NodeJS} and \textbf{ExpressJS} as the primary server-side language to develop the routes and connecting the server with database.
	
	\subsection{Database}
	We will be using \textbf{MongoDB} as the primary database to store all the necessary data of the Medicines and Appointments.
	
	\subsection{Integrated Development Environment:}
	\textbf{Visual Studio Code} is chosen as the suitable development environment to write, test, and debug our code efficiently.
	
	\subsection{Version Control}
	Using \textbf{Git} to track changes in our code, collaborate with team members, and maintain a version history.
	
	\subsection{Documentation Tools}
	\textbf{Google Docs} or \textbf {LaTeX}  used for project documentation.
	
	% ------------------------------------------------------------------------------
	
	% ------------------------------------------------------------------------------
	
	
	\section{Application Domain}
	
	The application domain for the IUT Medical Center Digital Platform includes patients seeking medical care, healthcare providers delivering services, administrative staff managing operations. \newline
	
	Patients require features such as appointment scheduling and access to medical records, while healthcare providers need tools for efficient care delivery. Administrative staff rely on the platform for tasks like room reservations and billing. By addressing the needs of these diverse user groups, the platform aims to enhance accessibility, efficiency, and overall user experience within the medical center.
	
	% ------------------------------------------------------------------------------
	
	% ------------------------------------------------------------------------------
	
	
	\section{Proposed Timeline}
	\begin{itemize}
		\item Week 1-2: Problem analysis and Requirements gathering
		\item Week 3-4: System Architecture and Proposal finalization
		\item Week 5-6: Organizing Data into database and Utility Methods
		\item Week 7-8: Development of Backend servers and routes
		\item Week 9-10: Implement the medical center features for routes
		\item Week 11-12: Development of Frontend client side UI
		\item Week 13: Integration and Testing
		\item Week 14: Bug fixes and Final Review
	\end{itemize}
	
	% ------------------------------------------------------------------------------
	
	% ------------------------------------------------------------------------------
	
\end{document}